%%%%%%%%%%%%%%%%%%%%%%%%%%%%%%%%%%%%%%%%%%%%%%%%%%%%%%%%%%%%%%%%%%%%%%%%%%%%%%%%
%%%%%%%%%%%%%%%%%%%%%%%%%%%%%%%%%%%%%%%%%%%%%%%%%%%%%%%%%%%%%%%%%%%%%%%%%%%%%%%%
%%%%%%%%%%%%%%%%%%%%%%%%%%%%%%%%%%%%%%%%%%%%%%%%%%%%%%%%%%%%%%%%%%%%%%%%%%%%%%%%
%%%%%%%%%%%%%%%%%%%%%%%%%%%%%%%%%%%%%%%%%%%%%%%%%%%%%%%%%%%%%%%%%%%%%%%%%%%%%%%%
\chapter[Correcteur PID]
{Différentes architectures des correcteurs PID\label{annexe-pid}}
%%%%%%%%%%%%%%%%%%%%%%%%%%%%%%%%%%%%%%%%%%%%%%%%%%%%%%%%%%%%%%%%%%%%%%%%%%%%%%%%
%%%%%%%%%%%%%%%%%%%%%%%%%%%%%%%%%%%%%%%%%%%%%%%%%%%%%%%%%%%%%%%%%%%%%%%%%%%%%%%%
%%%%%%%%%%%%%%%%%%%%%%%%%%%%%%%%%%%%%%%%%%%%%%%%%%%%%%%%%%%%%%%%%%%%%%%%%%%%%%%%
%%%%%%%%%%%%%%%%%%%%%%%%%%%%%%%%%%%%%%%%%%%%%%%%%%%%%%%%%%%%%%%%%%%%%%%%%%%%%%%%
%%%%%%%%%%%%%%%%%%%%%%%%%%%%%%%%%%%%%%%%%%%%%%%%%%%%%%%%%%%%%%%%%%%%%%%%%%%%%%%%
%%%%%%%%%%%%%%%%%%%%%%%%%%%%%%%%%%%%%%%%%%%%%%%%%%%%%%%%%%%%%%%%%%%%%%%%%%%%%%%%
\subsection*{PID idéal}
%%%%%%%%%%%%%%%%%%%%%%%%%%%%%%%%%%%%%%%%%%%%%%%%%%%%%%%%%%%%%%%%%%%%%%%%%%%%%%%%
%%%%%%%%%%%%%%%%%%%%%%%%%%%%%%%%%%%%%%%%%%%%%%%%%%%%%%%%%%%%%%%%%%%%%%%%%%%%%%%%
Pour élaborer un correcteur Proportionnel Intégral Dérivé (PID), il suffit
de construire un correcteur avec les propritétés des trois correcteurs 
élémentaires P, I et D comme brique de base. Il existe trois façon d'agencer
ces blocs de base pour élaborer un correcteur PID:
%-------------------------------------------------------------------------------
\begin{itemize}
    \item en \textbf{parallèle}
    \item en \textbf{série}
    \item et une combinaison des deux précédentes, dite \textbf{mixte}.
\end{itemize}
%-------------------------------------------------------------------------------
%%%%%%%%%%%%%%%%%%%%%%%%%%%%%%%%%%%%%%%%%%%%%%%%%%%%%%%%%%%%%%%%%%%%%%%%%%%%%%%%
\paragraph{parallèle}
%%%%%%%%%%%%%%%%%%%%%%%%%%%%%%%%%%%%%%%%%%%%%%%%%%%%%%%%%%%%%%%%%%%%%%%%%%%%%%%%
\index{Correcteur!Proportionnel Intégral Dérivé (PID)! parallèle}
Un PID en parallèle est la structure du correcteur la plus évidente, en effet
on combine directement la commande de
%-------------------------------------------------------------------------------
\begin{center}
    \includegraphics[width=0.7\textwidth]{pid_parallele_annexe_pid.pdf}
\end{center}
%-------------------------------------------------------------------------------
La réduction de ce schéma-bloc nous donne la relation entre la commande et 
l'écart
\[
    U(p)=\left(k_p^{(1)}+\dfrac{1}{\tau_i^{(1)} p}
                        +\tau_d^{(1)} p\right)\epsilon(p)
\]
%-------------------------------------------------------------------------------
\begin{bequation}[ams align]
    C_{\text{PID}}^{(1)}(p)=k_p^{(1)}+\dfrac{1}{\tau_i^{(1)} p}+\tau_d^{(1)} p
\end{bequation}
%-------------------------------------------------------------------------------
%%%%%%%%%%%%%%%%%%%%%%%%%%%%%%%%%%%%%%%%%%%%%%%%%%%%%%%%%%%%%%%%%%%%%%%%%%%%%%%%
\paragraph{série}
%%%%%%%%%%%%%%%%%%%%%%%%%%%%%%%%%%%%%%%%%%%%%%%%%%%%%%%%%%%%%%%%%%%%%%%%%%%%%%%%
\index{Correcteur! Proportionnel Intégral Dérivé (PID)! série}
%-------------------------------------------------------------------------------
\begin{center}
    \includegraphics[width=0.8\textwidth]{pid_serie_annexe_pid.pdf}
\end{center}
%-------------------------------------------------------------------------------
\[
    X(p)=k_p^{(2)}\left(1+\dfrac{1}{\tau_i^{(2)}p}\right)\epsilon(p)
\]
\[
    U(p)=X(p)\left(1+\tau_d^{(2)}p\right)
        =k_p^{(2)}\left(1+\dfrac{1}{\tau_i^{(2)}p}\right)
                  \left(1+\tau_d^{(2)}p\right)\epsilon(p)
\]
%-------------------------------------------------------------------------------
\begin{bequation}[ams align]
    C_{\text{PID}}^{(2)}(p)=k_p^{(2)}\left(1+\dfrac{1}{\tau_i^{(2)}p}\right)
                           \left(1+\tau_d^{(2)}p\right)
\end{bequation}
%-------------------------------------------------------------------------------
%%%%%%%%%%%%%%%%%%%%%%%%%%%%%%%%%%%%%%%%%%%%%%%%%%%%%%%%%%%%%%%%%%%%%%%%%%%%%%%%
\paragraph{mixte}
%%%%%%%%%%%%%%%%%%%%%%%%%%%%%%%%%%%%%%%%%%%%%%%%%%%%%%%%%%%%%%%%%%%%%%%%%%%%%%%%
\index{Correcteur!Proportionnel Intégral Dérivé (PID)! mixte}
%-------------------------------------------------------------------------------
\begin{center}
    \includegraphics[width=0.7\textwidth]{pid_mixte_annexe_pid.pdf}
\end{center}
%-------------------------------------------------------------------------------
\[
    U(p)=k_p^{(3)}\left(1+\dfrac{1}{\tau_i^{(3)}p}
                         +\tau_d^{(3)}p\right)\epsilon(p)
\]
%-------------------------------------------------------------------------------
\begin{bequation}[ams align]
    C_{\text{PID}}^{(3)}(p)=k_p^{(3)}\left(1+\dfrac{1}{\tau_i^{(3)}p}
                                            +\tau_d^{(3)}p\right)
\end{bequation}
%-------------------------------------------------------------------------------
Les trois structures du correcteur PID précédentes sont équivalentes, 
les paramètres $k_p,\tau_i,\tau_d$ n'auront simplement pas les mêmes valeurs.
Il est possible d'exprimer des relations de passage entre ces trois structures :
%%%%%%%%%%%%%%%%%%%%%%%%%%%%%%%%%%%%%%%%%%%%%%%%%%%%%%%%%%%%%%%%%%%%%%%%%%%%%%%%
\paragraph{3$\Rightarrow$ 1 et 1$\Rightarrow$ 3}
%%%%%%%%%%%%%%%%%%%%%%%%%%%%%%%%%%%%%%%%%%%%%%%%%%%%%%%%%%%%%%%%%%%%%%%%%%%%%%%%
\[
\begin{cases}
    k_p^{(1)}=k_p^{(3)}\\[1.0em]
    \tau_i^{(1)}=\dfrac{\tau_i^{(3)}}{k_p^{(3)}}\\[1.0em]
    \tau_d^{(1)}=k_p^{(3)}\tau_d^{(3)}
\end{cases}
\begin{cases}
    k_p^{(3)}=k_p^{(1)}\\[1.0em]
    \tau_i^{(3)}=k_p^{(1)}\tau_i^{(1)}\\[1.0em]
    \tau_d^{(3)}=\dfrac{\tau_d^{(1)}}{k_p^{(1)}}
\end{cases}
\]
%%%%%%%%%%%%%%%%%%%%%%%%%%%%%%%%%%%%%%%%%%%%%%%%%%%%%%%%%%%%%%%%%%%%%%%%%%%%%%%%
\paragraph{2$\Rightarrow$ 1 et 1$\Rightarrow$ 2}
%%%%%%%%%%%%%%%%%%%%%%%%%%%%%%%%%%%%%%%%%%%%%%%%%%%%%%%%%%%%%%%%%%%%%%%%%%%%%%%%
\[
\begin{cases}
    k_p^{(1)}=k_p^{(2)}\left(
              1+\dfrac{\tau_d^{(2)}}{\tau_i^{(2)}}\right)\\[1.0em]
    \tau_i^{(1)}=\dfrac{\tau_i^{(2)}}{k_p^{(2)}}\\[1.0em]
    \tau_d^{(1)}=k_p^{(2)}\tau_d^{(2)}
\end{cases}
\begin{cases}
k_p^{(2)}=
\dfrac{k_p^{(1)}}{2}\left(1+\sqrt{1-4\dfrac{\tau_d^{(1)}}{\tau_i^{(1)}(k_p^{(1)})^2}}\right)
\\[1.2em]

\tau_i^{(2)} = 
\dfrac{\tau_i^{(1)}k_p^{(1)}}{2}\left(1+\sqrt{1-4\dfrac{\tau_d^{(1)}}{\tau_i^{(1)}(k_p^{(1)})^2}}\right)\\[1.2em]

\tau_d^{(2)} = \dfrac{2\tau_d^{(1)}}{k_p^{(1)}\left(1+\sqrt{1-4\dfrac{\tau_d^{(1)}}{\tau_i^{(1)}(k_p^{(1)})^2}}\right)}
\end{cases}
\]
%%%%%%%%%%%%%%%%%%%%%%%%%%%%%%%%%%%%%%%%%%%%%%%%%%%%%%%%%%%%%%%%%%%%%%%%%%%%%%%%
\paragraph{2$\Rightarrow$ 3 et 3$\Rightarrow$ 2}
%%%%%%%%%%%%%%%%%%%%%%%%%%%%%%%%%%%%%%%%%%%%%%%%%%%%%%%%%%%%%%%%%%%%%%%%%%%%%%%%
\[
\begin{cases}
    k_p^{(3)}=k_p^{(2)}\left(
              1+\dfrac{\tau_d^{(2)}}{\tau_i^{(2)}}\right)\\[1.0em]
    \tau_i^{(3)}=\tau_i^{(2)}+\tau_d^{(2)}\\[1.0em]
    \tau_d^{(3)}=\dfrac{\tau_i^{(2)}\tau_d^{(2)}}{\tau_i^{(2)}+\tau_d^{(2)}}
\end{cases}
\begin{cases}
    k_p^{(2)}=k_p^{(3)}\left(
              1+\dfrac{\tau_d^{(3)}}{\tau_i^{(3)}}\right)\\[1.0em]
    \tau_i^{(2)}=\tau_i^{(3)}+\tau_d^{(3)}\\[1.0em]
    \tau_d^{(2)}=\dfrac{\tau_i^{(3)}\tau_d^{(3)}}{\tau_i^{(3)}+\tau_d^{(3)}}
\end{cases}
\]
